% =============================================================================
\chapter*{Introduction générale}
\addcontentsline{toc}{chapter}{Introduction générale}
\label{ch:introduction}
% =============================================================================

La virtualisation permet de mutualiser efficacement des charges hétérogènes sur un même serveur physique, mais cette mutualisation reste principalement intra-nœud. Dans un cluster Proxmox~VE, l’admission d’une machine virtuelle est d’abord contrainte par les ressources locales du nœud cible : un hôte proche de la saturation mémoire peut refuser une VM, même si la capacité agrégée du cluster demeure suffisante~\cite{proxmox_doc}. Ce découplage entre capacité globale et capacité locale conduit à une fragmentation opérationnelle : sous-utilisation de certains nœuds, pression forte sur d’autres, et migrations correctives plus fréquentes.

Le projet \textbf{FusionCore} (dépôt \texttt{omega-remote-paging}) traite ce problème en introduisant une couche de mutualisation \emph{en espace utilisateur}, sans modification du noyau Linux ni du code des VMs invitées. L’approche combine deux plans complémentaires :
\begin{itemize}
    \item un \textbf{plan de données local} en Rust (\texttt{omega-daemon}) pour les mécanismes sensibles à la latence (interception des défauts de page, éviction, service de pages distantes, collecte cgroups/PSI, interfaces QMP) ;
    \item un \textbf{plan de contrôle cluster} en Python (\texttt{omega-controller}) pour les décisions d’admission, de placement et de migration.
\end{itemize}

Le cœur du système est le \emph{paging distant}. Lorsqu’une page froide d’une VM est éjectée vers un autre nœud, la page reste adressable de manière transparente ; un futur accès déclenche un défaut de page géré via \texttt{userfaultfd}, puis une récupération réseau de la page et sa réinjection locale. L’architecture s’appuie sur des interfaces standards du noyau et de l’hyperviseur : \texttt{userfaultfd}, cgroups~v2, PSI, QMP, API Proxmox et stockage partagé Ceph~RBD~\cite{kvm_doc,qemu_doc,ceph_doc}.

Au-delà de la mémoire, FusionCore met en cohérence plusieurs mécanismes réellement présents dans le code source :
\begin{enumerate}
    \item quota mémoire par VM (budget local + distant) appliqué sur les écritures de pages ;
    \item ordonnancement vCPU local (hotplug/hot-unplug et ajustements cgroups) ;
    \item exposition et réservation VRAM via API de contrôle ;
    \item surveillance de la pression I/O et ajustement des poids \texttt{io.weight} ;
    \item migration live/cold déclenchée sur critères de pression et de faisabilité.
\end{enumerate}

\textbf{Problématique.} Dans quelle mesure peut-on améliorer l’utilisation effective d’un cluster Proxmox hétérogène en coordonnant mémoire, CPU et GPU sans dépendances matérielles spécialisées (RDMA/CXL) et sans altérer la chaîne de maintenance standard de l’infrastructure ?

\textbf{Objectif du mémoire.} Documenter de façon rigoureuse la conception et l’implémentation de FusionCore, en explicitant les choix techniques, les compromis, les flux de données et les limites observées dans l’environnement d’expérimentation du projet.

\textbf{Contributions du rapport.} Ce document propose :
\begin{itemize}
    \item une analyse du besoin et un positionnement par rapport aux travaux de désagrégation ;
    \item une description d’architecture fidèle aux composants implémentés ;
    \item une formalisation des invariants opérationnels (admission, quotas, migration) ;
    \item une méthodologie d’évaluation reproductible basée sur les scripts et scénarios du dépôt ;
    \item une discussion critique des limites actuelles et des pistes d’évolution.
\end{itemize}

Le manuscrit est structuré comme suit : le chapitre~\ref{ch:etat-art} présente l’état de l’art et positionne FusionCore ; le chapitre~\ref{cp:analyse-conception} formalise l’architecture et les choix de conception ; le chapitre~\ref{ch:implementation-resultats} décrit l’implémentation et la méthodologie expérimentale reproductible ; enfin le chapitre~\ref{ch:sma} interprète la coordination globale sous l’angle des systèmes multi-agents.
